% =============================================================================
% Section 6 results fragment: the Economic Opportunity pillar (final run).
% Numbers: paper/results/eo.json; figures, numbers and the two tables built by
% scripts/paper/fig_eo.R from data/processed/outputs/index_outputs.rds (final rerun of
% 2026-10-12). Every cross-country statistic (distribution, fixed effects, contributions, ranks,
% leaders, coalition positions, top tens) uses rankable cells only (config/rankability.yml:
% sovereign economy and at most half of the cell's index weight imputed); rank correlations are
% Spearman.
% Figures (file -> label): eo_distribution -> fig:eo_distribution, eo_contrib -> fig:eo_contrib,
% eo_heatmap -> fig:eo_heatmap, eo_top -> fig:eo_top. This fragment replaces the four \tbd
% paragraphs and the four placeholder floats of subsec:eo_results (the placeholder label
% fig:eo_category_contrib becomes fig:eo_contrib). Tables: tables/eo_weight_split.tex
% (tab:eo_weight_split) and tables/eo_leaders.tex (tab:eo_leaders), written by fig_eo.R.
% =============================================================================

\paragraph{Coverage and rankability.} Among the eleven headline technologies the pillar is scored
for 7{,}131 country-technology-stage cells, which cover 219 economies (195 sovereign economies and
24 territories) and all 33 technology-stage pairs; each pair scores between 213 and 219 economies.
All 33 pairs have status \emph{ok} in the diagnostics of Section~\ref{subsec:norm_diagnostics}: none
is constant across economies and none is unscored. Within each pair every economy is scored on the
same set of categories, so a comparison within a pair is a comparison of one construct; the number
of categories per cell is between three and seven upstream, three and six midstream and two and
seven downstream. Twenty-four cells rest entirely on imputed components ($\iota_{its}=1$ in
equation~\eqref{eq:imputed_share}), all of them in two territories and one microstate: the Isle of
Man (14 cells), Puerto Rico (8) and Liechtenstein (2). Short of that extreme, imputation is
pervasive: the mean imputed weight share is 0.43 (0.41 upstream, 0.44 midstream and 0.43
downstream), and 45.3 per cent of the category scores contain at least one imputed component.
Because each component is taken at the latest year recorded for it
(Section~\ref{subsec:norm_latest}), 92.1 per cent of the scored cells combine index components from
more than one year, and a cell combines up to six. Both decompositions of
Section~\ref{subsec:norm_decomposition} reconcile: the largest absolute gap between $\EO_{its}$ and
the sum of its category contributions is zero, and between $\EO_{its}$ and the sum of its variable
contributions it is $1.1\times10^{-16}$.

The rankability rule of Section~\ref{subsec:norm_diagnostics} admits 4{,}301 of the 7{,}131 cells
(60.3 per cent), for 178 economies, to the cross-country comparisons that follow. Of the 2{,}830
cells it sets aside, 2{,}112 carry more than half of their weight on imputed components, 74 belong
to territories, and 644 fail on both counts. The rule bites hardest where the pillar's
country-specific sources are thinnest: 70.2 per cent of upstream cells are rankable, against 55.7
per cent midstream and 54.9 per cent downstream, and the number of rankable economies per pair
ranges from 40 (green hydrogen midstream) and 54 (batteries midstream) to 175 (nuclear upstream),
with a median of 142. In the clean-technology midstream pairs the rule leaves mainly the economies with recorded manufacturing capacity or investment: in batteries 42 of the 54
rankable economies have an observed Foreign Dependency or Investment input, in green hydrogen 30
of 40. The mean imputed weight share of the rankable cells is 0.29 (0.32 upstream, 0.28 midstream
and 0.25 downstream). The excluded cells would mostly have ranked low. Ranked among all scored
economies of their pair, the cells excluded for imputation alone have a mean percentile of 0.43
and 63.4 per cent of them lie below the pair median; territory cells have 0.42 and 63.8 per cent,
against 0.55 and 41.9 per cent for the rankable cells. Of the 330 top-ten places computed over all
scored economies, only three, and none of the first places, belong to cells the rule excludes. The
rule therefore changes the composition of the lower half of each pair far more than its top.

Technological Readiness is mapped to stages: each item of the IEA Clean Technology Guide is
assigned to a technology and a stage, and a stage without items takes the technology-wide value.
The TRL Index takes the same value for every economy in a pair, but it differs across the three
stages of ten of the eleven technologies (wind keeps one value), with component values from 0.31
(geothermal downstream) to 0.72 (gas downstream) and a median range across a technology's stages
of 0.16. It therefore moves the level of a cell relative to the other stages of the same
technology, which matters for the within-country Strategic Index, but it cannot reorder economies
within a pair.

\paragraph{Distribution.} Figure~\ref{fig:eo_distribution} shows the pillar by technology within
each stage, over the rankable economies of each pair. The stages are alike in level: the median
rankable cell scores 0.500 upstream, 0.501 midstream and 0.488 downstream (means 0.487, 0.489 and
0.489), with pooled interquartile ranges of 0.434 to 0.542, 0.427 to 0.561 and 0.427 to 0.551.
The similarity is partly the rule's doing. Over all scored cells the midstream median is 0.454 and
its lower quartile 0.328, because the zero rule places the many economies with no recorded
manufacturing capacity or investment at the bottom of the clean-technology midstream pairs, and
those economies, which carry most of their weight on filled inputs, are the ones the rule removes.
The figure is dominated by differences between technologies rather than between economies. Pair
medians range from 0.367 (coal) to 0.585 (electric grid) upstream, from 0.351 (wind) to 0.579
(oil) midstream and from 0.363 (geothermal) to 0.546 (oil) downstream, whereas the mean
interquartile range across economies within a pair is 0.037 upstream, 0.082 midstream and 0.072
downstream. Four of the five clean-technology midstream pairs, wind, electric vehicles, solar and
batteries, have the lowest midstream medians (0.351 to 0.426); green hydrogen (0.460) sits next to
coal (0.457), and the oil, geothermal, gas, nuclear and electric-grid pairs lie between 0.50 and
0.58. The ordering is a construction effect rather than a finding about the technologies. The four
low pairs carry Foreign Dependency and Investment under the zero rule, so even among the rankable
economies those without recorded capacity or investment score zero on up to two categories,
whereas the other midstream pairs are scored on Trade, Technology Demand, Technological Readiness
and Cost Competitiveness, none of which carries the zero rule.

The upstream boxes are the narrowest because the upstream pairs carry the largest weight on inputs
that do not vary across economies (Table~\ref{tab:eo_weight_split}). The widest boxes owe their
width to one category each. In electric vehicles midstream (interquartile range 0.144), Foreign
Dependency, the electric-vehicle production position, is observed for 97 of the 145 rankable
economies and zero for the rest; it accounts for 46 per cent of the pair's variance, and with it
held at its pair mean the interquartile range would be 0.053. In green hydrogen downstream (0.122)
the Consumption category accounts for 94 per cent of the variance (interquartile range 0.017
without it), and in wind downstream (0.117) Investment, observed for 61 of the 95 rankable
economies, accounts for 46 per cent (0.063 without it).

\paragraph{Variance decomposition.} Pooled over the 4{,}301 rankable cells, technology-stage fixed
effects explain 57.9 per cent of the variance of $\EO_{its}$. Technology effects alone explain 27.8
per cent and stage effects alone almost nothing (0.01 per cent), so the level of a pair is set by
its specific composition, which categories it carries and under which missing-data rules, rather
than by the stage as such. Country fixed effects alone explain 17.7 per cent; together with the
pair effects they explain 77.2 per cent, so country effects account for 45.8 per cent of the
variance that remains within pairs. On the same definition the energy-security pillar
(Section~\ref{subsec:es_results}), over its 4{,}163 rankable cells in 30 pairs, has a pair share of
44.4 per cent and a within-pair country share of 27.8 per cent. Over all 7{,}131 scored cells the
pair share of $\EO$ is 71.5 per cent, higher in part because the zero-filled cells that the rule
removes pull the clean-technology midstream pairs down. The large pair share is anticipated by
construction. Technological Readiness, Energy Prices and most of Technology Demand are constants of
the pair, and by equation~\eqref{eq:affine_constant} such a category shifts the level of a pair
and compresses the spread within it without reordering its economies. Comparisons of $\EO$ levels
across pairs, such as those the within-country Strategic Index makes (Section~\ref{sec:strategic}),
therefore compare mostly pair effects; what the pillar measures about economies is chiefly their
order within a pair.

That order is persistent across the supply chain, and the large within-pair country share has a
common source. Across the 159 economies rankable at all three stages, the Spearman rank correlation
between an economy's mean within-pair percentile rank (equation~\eqref{eq:pct_rank}, ties averaged)
at two stages is 0.86 for upstream and midstream, 0.70 for upstream and downstream and 0.79 for
midstream and downstream. Part of this common factor is mechanical. The Input Cost Index, the one
component of Cost Competitiveness defined in every technology and stage, is a country-level index
of wages and capital costs whose stage variation comes only from the activity used for wages and
from the factor shares of Table~\ref{tab:eo_factor_shares}. Across the 178 economies with a
rankable cell, the Spearman correlation between an economy's mean within-pair percentile on $\EO$
and its mean within-pair percentile on the Cost Competitiveness category is 0.77, and 0.82 across
the 57 coalition members.

\paragraph{What drives the pillar.} Figure~\ref{fig:eo_contrib} reports, by stage, the mean share
of each category in the level of the pillar, from the category decomposition of
equation~\eqref{eq:cat_decomp}, and its mean share of the variance of $\EO_{its}$ across the
rankable economies of a pair, $\mathrm{cov}(\omega_{kits}c_{kits},\EO_{its})/\mathrm{var}(\EO_{its})$,
which sums to one over categories. The two panels tell different stories. In the level, the
largest shares are Cost Competitiveness (21.9 per cent), Technology Demand (17.3), Production
(16.5) and Energy Prices (16.2) upstream; Cost Competitiveness (29.2), Technology Demand (24.7) and
Trade (24.6) midstream; and Cost Competitiveness (27.5), Technology Demand (19.7) and Technological
Readiness (18.0) downstream. The three categories that are wholly or largely constant within a
pair, Technology Demand, Technological Readiness and Energy Prices, carry a combined mean effective
weight (equation~\eqref{eq:effective_weight}) of 48.5 per cent upstream, 42.2 per cent midstream
and 35.2 per cent downstream. In the variance across economies, Technological Readiness and Energy
Prices contribute exactly nothing and Technology Demand between 1.9 and 5.7 per cent, its only
country-varying components being the BNEF regional composite, observed for at most twelve
economies, and the EV-market composite. The categories that order economies are Cost
Competitiveness (39.2 per cent of the within-pair variance upstream, 37.6 midstream and 21.7
downstream), Trade (26.9) and Production (22.4) upstream, Trade (25.5) and Foreign Dependency
(22.2) midstream, and Consumption (23.8) and Investment (22.2) downstream. Upstream, Trade
orders economies in the nine pairs with a trade basket, among them the electric-vehicle, wind,
solar and nuclear pairs, whose baskets of battery-mineral ores, rare earths, silicon metal and
uranium the final basket review added. Investment illustrates the difference between the two
measures: it contributes 0.3, 1.2 and 4.2 per cent of the level upstream, midstream and
downstream, but 8.3, 12.7 and 22.2 per cent of the within-pair variance, because under the zero
rule the economies with recorded investment stand far above those with none, which are the
majority of the rankable economies in 9 of the 12 pairs that carry the category.

% Written by scripts/paper/fig_eo.R from the objects saved in paper/results/eo.json (keys: input_weights, coverage).
\begin{table}[tbp]
\centering
\begin{threeparttable}
\caption{Effective weight of the Economic Opportunity pillar by type of input}
\label{tab:eo_weight_split}
\small
\begin{tabular}{@{}l S[table-format=2.1] S[table-format=2.1] S[table-format=2.1] S[table-format=2.1]@{}}
\toprule
 & \multicolumn{3}{c}{Mean effective weight (per cent)} & {Constants' share}\\
\cmidrule(lr){2-4}
 & {Technology-level} & {Imputed,} & {Observed,} & {of the level}\\
 & {constants} & {country-varying} & {country-varying} & {(per cent)}\\
\midrule
Upstream & 46.3 & 19.7 & 33.9 & 44.0\\
Midstream & 40.4 & 14.6 & 45.0 & 40.2\\
Downstream & 27.9 & 17.4 & 54.7 & 30.4\\
\addlinespace
All rankable cells & 38.9 & 17.5 & 43.6 & 38.7\\
All scored cells & 37.7 & 30.1 & 32.3 & 39.5\\
\bottomrule
\end{tabular}
\begin{tablenotes}[flushleft]\footnotesize
\item \textit{Notes:} The effective weight of variable $v$ in cell $(i,t,s)$ is the renormalised category weight times the variable's weight within its category, $\omega_{kits}\lambda_{vkits}$ (equation~\eqref{eq:var_decomp}); it sums to one in every cell. Technology-level constants are the variables whose value has zero range across all scored economies of the pair: the TRL Index (one value per technology and stage), the WEO global-demand composite, the BCG addressable market, the overcapacity index, the world-market concentration term of Trade and commodity-price volatility. Imputed components are those filled by a missing-data rule (global average or zero) for the economy in question. The first three columns sum to 100 within a row, up to rounding. The last column is the mean share of $\EO_{its}$ contributed by the constants. Rows by stage and ``All rankable cells'': means over the 1{,}681 upstream, 1{,}307 midstream and 1{,}313 downstream rankable cells of the eleven headline technologies (a sovereign economy with at most half of the cell's weight imputed); ``All scored cells'': means over all 7{,}131 scored cells. Source: \texttt{economic\_opportunity\_variable\_contributions} in \texttt{data/processed/outputs/index\_outputs.rds}; \texttt{scripts/paper/fig\_eo.R}.
\end{tablenotes}
\end{threeparttable}
\end{table}


Table~\ref{tab:eo_weight_split} makes the point at the level of variables. Six variables are
constant across economies in every pair in which they appear: the stage-mapped TRL Index, the WEO
global-demand composite, the BCG addressable market (recorded for the United States and copied to
every other economy by the global-average rule), the overcapacity index, the world-market
concentration term of Trade (ranked across the cells of the eleven headline technologies only) and
commodity-price volatility. Together they carry 38.9 per cent of the effective weight of the
average rankable cell and contribute 38.7 per cent of its level; their weight is 46.3 per cent
upstream, where price volatility (15.9), the TRL Index (12.5) and the addressable market (10.1)
are the largest, 40.4 per cent midstream and 27.9 per cent downstream. Country-varying components
that are imputed for the economy in question carry a further 17.5 per cent, so that in the average
rankable cell 43.6 per cent of the weight rests on country-specific information observed for that
economy, from 33.9 per cent upstream to 54.7 per cent downstream. Over all scored cells, before the
rankability rule, the observed country-specific share is only 32.3 per cent and the imputed share
30.1 per cent. Even after the rule, the pillar's cross-country content is thinner than its ten
categories suggest, and it is thinnest upstream.

% Written by scripts/paper/fig_eo.R from the objects saved in paper/results/eo.json (key: leaders).
\begin{table}[tbp]
\centering
\begin{threeparttable}
\caption{Leading economies in each technology-stage pair of the Economic Opportunity pillar (rankable cells)}
\label{tab:eo_leaders}
\footnotesize
\begin{tabular}{@{}lrlllcrr@{}}
\toprule
Technology & $n$ & First & Second & Third & Best member (rank) & USA & CHN \\
\midrule
\multicolumn{8}{@{}l}{\textit{Upstream}}\\
Batteries & 120 & CHN 0.65 & IDN 0.59 & AUS 0.57 & IDN (2) & 7 & 1 \\
Coal & 157 & IDN 0.52 & CHN 0.51 & IND 0.48 & IDN (1) & 7 & 2 \\
Electric Grid & 163 & CHN 0.71 & CHL 0.65 & RUS 0.65 & CHL (2) & 27 & 1 \\
Electric Vehicles & 153 & CHN 0.68 & IDN 0.62 & AUS 0.57 & IDN (2) & 4 & 1 \\
Gas & 165 & USA 0.61 & CHN 0.60 & CAN 0.57 & USA (1) & 1 & 2 \\
Geothermal & 141 & TCD 0.48 & CMR 0.48 & COG 0.48 & PRT (4) & 68 & 6 \\
Green Hydrogen & 146 & CHN 0.62 & IDN 0.62 & ZMB 0.57 & IDN (2) & 62 & 1 \\
Nuclear & 175 & CAN 0.65 & CHN 0.62 & KAZ 0.62 & CAN (1) & 6 & 2 \\
Oil & 166 & USA 0.66 & CHN 0.64 & CAN 0.64 & USA (1) & 1 & 2 \\
Solar & 153 & CHN 0.70 & IDN 0.66 & THA 0.61 & IDN (2) & 23 & 1 \\
Wind & 142 & CHN 0.70 & IDN 0.59 & USA 0.58 & IDN (2) & 3 & 1 \\
\addlinespace
\multicolumn{8}{@{}l}{\textit{Midstream}}\\
Batteries & 54 & CHN 0.77 & USA 0.63 & HUN 0.55 & USA (2) & 2 & 1 \\
Coal & 155 & CHN 0.59 & POL 0.57 & IDN 0.56 & POL (2) & 14 & 1 \\
Electric Grid & 162 & CHN 0.68 & THA 0.64 & POL 0.62 & THA (2) & 14 & 1 \\
Electric Vehicles & 145 & CHN 0.74 & ESP 0.64 & USA 0.63 & ESP (2) & 3 & 1 \\
Gas & 156 & CHN 0.71 & IND 0.67 & ITA 0.65 & IND (2) & 5 & 1 \\
Geothermal & 141 & TCD 0.64 & CMR 0.64 & COG 0.64 & PRT (4) & 58 & 5 \\
Green Hydrogen & 40 & CHN 0.71 & DEU 0.50 & USA 0.50 & DEU (2) & 3 & 1 \\
Nuclear & 154 & CHN 0.65 & FRA 0.65 & IND 0.62 & FRA (2) & 21 & 1 \\
Oil & 156 & CHN 0.72 & AUT 0.67 & ITA 0.66 & AUT (2) & 11 & 1 \\
Solar & 79 & CHN 0.79 & IND 0.64 & USA 0.63 & IND (2) & 3 & 1 \\
Wind & 65 & CHN 0.74 & POL 0.54 & IND 0.52 & POL (2) & 4 & 1 \\
\addlinespace
\multicolumn{8}{@{}l}{\textit{Downstream}}\\
Batteries & 140 & CHN 0.75 & USA 0.67 & ITA 0.63 & USA (2) & 2 & 1 \\
Coal & 104 & CHN 0.68 & IDN 0.66 & IND 0.66 & IDN (2) & 15 & 1 \\
Electric Grid & 141 & TCD 0.68 & CMR 0.68 & TGO 0.67 & THA (11) & 109 & 20 \\
Electric Vehicles & 105 & CHN 0.75 & VNM 0.73 & THA 0.70 & VNM (2) & 29 & 1 \\
Gas & 104 & USA 0.72 & SAU 0.71 & CHN 0.68 & USA (1) & 1 & 3 \\
Geothermal & 128 & USA 0.61 & IDN 0.46 & PRT 0.42 & USA (1) & 1 & 16 \\
Green Hydrogen & 158 & USA 0.65 & DEU 0.60 & THA 0.60 & USA (1) & 1 & 4 \\
Nuclear & 86 & CHN 0.72 & FRA 0.67 & JPN 0.64 & FRA (2) & 15 & 1 \\
Oil & 104 & SAU 0.71 & IRN 0.69 & QAT 0.68 & SAU (1) & 8 & 98 \\
Solar & 148 & CHN 0.75 & USA 0.67 & IND 0.67 & USA (2) & 2 & 1 \\
Wind & 95 & CHN 0.70 & USA 0.61 & CAN 0.60 & USA (2) & 2 & 1 \\
\bottomrule
\end{tabular}
\begin{tablenotes}[flushleft]\footnotesize
\item \textit{Notes:} $\EO_{its}$ for the eleven headline technologies. Only rankable cells enter: a sovereign economy with at most half of the cell's index weight imputed (\texttt{config/rankability.yml}). Economies are identified by ISO3 code and followed by their score; $n$ is the number of rankable economies in the pair (40 to 175). ``Best member'' is the highest-ranked of the 57 coalition members (31 core, 26 partners), with its rank in parentheses; the USA and CHN columns give the ranks of the United States and China. A rank is one plus the number of rankable economies with a strictly higher score. All 33 pairs are scored and none is constant across economies. In geothermal upstream and midstream and electric-grid downstream, Cost Competitiveness is the only category that varies across economies. Codes: AUS Australia, AUT Austria, CAN Canada, CHL Chile, CMR Cameroon, COG Congo, DEU Germany, ESP Spain, FRA France, HUN Hungary, IDN Indonesia, IND India, IRN Iran, ITA Italy, JPN Japan, KAZ Kazakhstan, POL Poland, PRT Portugal, QAT Qatar, RUS Russia, SAU Saudi Arabia, TCD Chad, TGO Togo, THA Thailand, VNM Viet Nam, ZMB Zambia, USA United States, CHN China. Source: \texttt{economic\_opportunity\_index} in \texttt{data/processed/outputs/index\_outputs.rds}; \texttt{scripts/paper/fig\_eo.R}.
\end{tablenotes}
\end{threeparttable}
\end{table}


\paragraph{Leaders.} Table~\ref{tab:eo_leaders} lists the three highest-scoring rankable economies
in each pair and the best-placed coalition member. China leads 22 of the 33 pairs, the United
States five, Chad three, and Canada, Indonesia and Saudi Arabia one each; there are no ties at the
top. Classifying as emerging economies those the World Bank places in the low, lower-middle or
upper-middle income groups, 26 of the 33 leaders are emerging economies, but the count is almost
entirely China's. The four other emerging leaders are Indonesia in coal upstream and Chad in
geothermal upstream, geothermal midstream and electric-grid downstream. Chad's three pairs are
precisely those in which Cost Competitiveness is the only category that varies across economies,
so that $\EO$ there is a ranking of the Input Cost Index. Chad combines low recorded wages (a mean
percentile of 0.12 among the economies reporting earnings) with a recorded lending rate of 2.45 per
cent, a 2017 observation that is the joint fifth lowest of the 144 rates in the cost theme and sits
at the floor of the rescaled rate $R_i$ of equation~\eqref{eq:eo_capital}. The runner-up in all
three pairs is Cameroon, which records the same rate (in 2018), so the two economies differ only in
wages and producer-price inflation; Chad's margins over it are 0.0020, 0.0005 and 0.0011. These
three leads are artefacts of pairs with almost no country-specific content and should not be read
as evidence of opportunity. Canada leads nuclear upstream (0.65, ahead of China and Kazakhstan),
and China leads electric-grid upstream (0.71) ahead of Chile (0.65), the pair whose trade category
is built on the copper and aluminium input basket. Margins at the top are generally small: the
median gap between the leader and the second economy is 0.043 (0.058 where China leads), and it is
below 0.01 in seven pairs. The largest gaps are China's in green hydrogen midstream (0.202 over
Germany), wind midstream (0.201 over Poland), solar midstream (0.149 over India) and batteries
midstream (0.138 over the United States), together with that of the United States in geothermal
downstream (0.153 over Indonesia). Below the first place the picture is more mixed: across the top
ten of all 33 pairs, emerging economies hold 46.4 per cent of the 330 places (41.0 per cent of the
places not held by China) and coalition members 69.1 per cent; China is in the top ten of 30 pairs
and the United States of 20.

The United States ranks in the bottom half of only one pair, and the reason is again
construction: it is 109th of 141 in electric-grid downstream, a pair in which Cost
Competitiveness is the only category that varies across economies and in which its Input Cost
Index sits at a within-pair percentile of 0.23. In the mineral-based upstream pairs, where mined
production is scored as a demand-weighted share of world mined supply and its growth as an
absolute change, it ranks between third (wind) and 62nd (green hydrogen). China's one low rank, 98th of 104
in oil downstream, comes from its Consumption category, on which it ranks last in the pair, and
from the growth component of its Production category (within-pair percentile 0.04).

\paragraph{The coalition.} Figure~\ref{fig:eo_heatmap} shows, for the 31 core members and 26
partners of the coalition (Section~\ref{sec:coalition}) and for China, each economy's percentile
rank within each pair among the pair's rankable economies. The coalition stands above the rest of
the world but well below China, and its two tiers differ. Over their rankable cells, the core
members have a mean within-pair percentile of 0.66 and a mean $\EO$ of 0.510, the partners 0.49
and 0.485, the 120 other rankable economies 0.44 and 0.479, and China 0.96 and 0.659. The core is
strongest midstream (mean percentile 0.73, against 0.64 upstream and 0.59 downstream), whereas the
partners stay near the middle of their pairs at every stage (0.48, 0.47 and 0.51). In most pairs
some member ranks high: the best-placed member is the world leader in 8 pairs and in the top three
in 30, with a median rank of 2. The three exceptions are the cost-only pairs that Chad leads, where
the best member ranks fourth (Portugal, in both geothermal pairs) and eleventh (Thailand,
electric-grid downstream). In 14 pairs the best-placed member is a partner (Indonesia in seven,
India and Thailand in two each), and upstream the partners supply the strongest member candidate
in seven of the eleven pairs: Indonesia ranks second behind China in five mineral pairs and first
in coal, and Chile second in electric grid, while the best core member ranks between third and
seventh in these pairs.

Among the core members, the United States (mean percentile 0.89), Portugal (0.84), Poland (0.84),
the Czech Republic (0.81) and Spain (0.80) rank highest, and the United Kingdom (0.57), Romania
(0.52), Australia (0.46), New Zealand (0.34) and Ukraine (0.11) lowest; among the partners,
Thailand (0.82), Indonesia (0.80), Malaysia (0.79), India (0.78) and Saudi Arabia (0.66) rank
highest, and the Dominican Republic (0.26), Brazil (0.23), Argentina (0.19), the Democratic
Republic of the Congo (0.13) and Ghana (0.10) lowest. Both ends of the list are largely a cost
effect. Four members, T\"urkiye, Brazil, Argentina and Ghana, record nominal lending rates (39.0,
46.0, 59.7 and 27.0 per cent, the 141st, 142nd, 143rd and 139th of the 144 rates in the cost
theme) that sit at the winsorisation ceiling of the rescaled rate $R_i$
(equation~\eqref{eq:eo_capital}); since the Input Cost Index enters every pair, the rate depresses
each of them across the board. With the Cost Competitiveness contribution replaced by its pair mean
over the rankable economies, T\"urkiye's mean percentile would be 0.73 rather than 0.30, Brazil's
0.61 rather than 0.23, Argentina's 0.45 rather than 0.19 and Ghana's 0.39 rather than 0.10;
Ukraine, whose recorded rate of 18.9 per cent is the 130th of 144, would move from 0.11 to 0.50.
Six members gain more than 0.2 in this way. At the top the same category works in the other
direction: neutralising it would lower Portugal from 0.84 to 0.62, Thailand from 0.82 to 0.64 and
Ireland from 0.71 to 0.51, while the United States (0.89 to 0.87) barely moves. A second
construction effect comes from the BNEF regional composite of Technology Demand, which is observed
for at most twelve economies; the global-average rule gives every other economy their mean, so an
economy that BNEF models but whose value lies below that mean ranks below the many economies it
does not model. Australia's value lies below the mean in all 18 pairs in which it is observed, the
United Kingdom's in 14 of 18 and France's in 15 of 21; with Technology Demand at its pair mean,
Australia's mean percentile would be 0.52 rather than 0.46, the United Kingdom's 0.65 rather than
0.57 and France's 0.80 rather than 0.76, while Viet Nam and India, whose modelled values lie mostly
above the mean, would fall from 0.63 to 0.57 and from 0.78 to 0.76. Readers who regard nominal
rates in high-inflation economies or market-rate wage levels as poor guides to the cost of new
capacity, or who doubt that an economy modelled by BNEF should rank below those it does not model,
should discount these positions accordingly.

Coverage differs between the tiers. Of the core members' 1{,}023 scored cells 1{,}002 are
rankable, against 716 of the partners' 858; the cells set aside are member cells with more than
half of their weight imputed, shown as dotted gray cells in Figure~\ref{fig:eo_heatmap}. Gabon has
only 9 rankable cells, Oman 21, Saudi Arabia and Singapore 22 each and Morocco, Namibia and Costa
Rica 23, so their mean positions rest on a minority or a bare majority of the pairs. Among the
rankable cells the mean imputed weight share is 0.20 for the core, 0.27 for the partners, 0.33 for
the other economies and 0.10 for China, and it is zero for the United States, partly because the
BCG addressable market is recorded for the United States alone. Imputation and rank remain related
within pairs: the median, across the 33 pairs, of the within-pair Spearman correlation between
$\EO_{its}$ and $\iota_{its}$ over the rankable economies is $-0.28$ (range $-0.94$ to $0.17$).
Under the global-average rule an uncovered economy is placed at the mean of the covered ones and
under the zero rule at the bottom, so the more heavily imputed of the rankable economies still tend
to sit lower in their pairs.

\paragraph{Illustrative pairs.} Figure~\ref{fig:eo_top} decomposes the scores of the ten
highest-ranked economies in upstream oil, midstream batteries, midstream solar and downstream
electric vehicles. The first three are the pairs of Figure~\ref{fig:es_top}; upstream oil is,
after nuclear upstream (175 economies), the upstream pair with the most rankable economies (166).
In each panel the gray segment, the sum of the technology-level constants, has the same length for
every economy (0.250 in oil, 0.202 in batteries, 0.181 in solar and 0.088 in electric vehicles),
so the order is set entirely by the coloured segments. Upstream oil is led by the United States
(0.659), China (0.644), Canada (0.637), Saudi Arabia (0.622) and Norway (0.607), with four
coalition members in the top ten. The differences among them come from Production, Trade and Cost
Competitiveness, whose contributions span 0.051, 0.039 and 0.027 across the ten: the United States
has the largest Production (0.167) and Trade (0.114) contributions, China the largest Cost
Competitiveness contribution (0.122). Saudi Arabia, Iraq, Iran, Kuwait and the United Arab
Emirates, five of the ten, have the same Cost Competitiveness contribution, 0.095, the smallest in
the panel and equal to its largest imputed country-varying contribution, and they carry between 32
and 38 per cent of their weight on imputed components; their places rest on production and trade.
In midstream batteries, where only 54 economies are rankable, China (0.771) leads the United States
(0.633) by 0.138, a gap that comes from Foreign Dependency (0.167 against 0.106, China's share of
world manufacturing capacity) and Cost Competitiveness (0.142 against 0.073); the Investment
contribution of the United States (0.143) exceeds China's (0.134). Nine of the ten are coalition
members, six of them core members. Midstream solar shows the same structure: China (0.788) is
followed at a distance by India (0.640) and the United States (0.635), five Southeast Asian
partners, Viet Nam, Thailand, Indonesia, Malaysia and the Philippines, take five of the remaining
seven places, and France and Germany are eighth and ninth. Emerging economies are most numerous at
the top of downstream electric vehicles, where they take eight of the ten places (against seven in
midstream solar, four in midstream batteries and three in upstream oil): China (0.755), Viet Nam
(0.731), Thailand (0.705) and T\"urkiye (0.702) lead. The pair's largest segment is the EV-market
composite of Technology Demand, which contributes between 0.21 and 0.33 to the top ten and half of
whose twelve indices are growth rates of EV stock, sales and sales share
(equation~\eqref{eq:eo_ev}). The composite therefore rewards markets that are growing fast from a
small base as well as large ones: T\"urkiye, Viet Nam and the Philippines receive larger Technology
Demand contributions than China (0.333, 0.295 and 0.281 against 0.271), so the pair measures market
momentum at least as much as market scale. T\"urkiye reaches fourth place despite the smallest Cost
Competitiveness contribution of the ten (0.057), the mark of its recorded lending rate.

\begin{figure}[tbp]
  \centering
  \includegraphics[width=\textwidth]{eo_distribution}
  \caption{Distribution of the Economic Opportunity pillar by technology and supply-chain stage.
  Each panel is one stage; each box summarises $\EO_{its}$ across the rankable economies of one
  technology-stage pair (their number in parentheses, 40 to 175), showing the median and the
  interquartile range, with whiskers to the most extreme economy within 1.5 times the
  interquartile range of the box and the remaining economies as dots. A cell is rankable if the
  economy is sovereign and at most half of the cell's index weight is imputed
  (\texttt{config/rankability.yml}). Technologies are ordered by their median within each panel.
  The score lies in $[0,1]$, and higher values mean more opportunity. All 33 pairs have status
  \emph{ok} (none is constant across economies or unscored). Sample: the 4{,}301 rankable cells
  of the eleven headline technologies (of 7{,}131 scored). Source:
  \texttt{economic\_opportunity\_index} in \texttt{data/processed/outputs/index\_outputs.rds};
  \texttt{scripts/paper/fig\_eo.R}.}
  \label{fig:eo_distribution}
\end{figure}

\begin{figure}[tbp]
  \centering
  \includegraphics[width=\textwidth]{eo_contrib}
  \caption{What the Economic Opportunity pillar is made of, by supply-chain stage. Left: the mean
  share of each category in the level of the pillar, $\omega_{kits}c_{kits}/\EO_{its}$ from the
  category decomposition of equation~\eqref{eq:cat_decomp}, averaged over the rankable cells of the
  stage (a category absent from a cell counts as zero). Right: the mean share of each category in
  the variance of the pillar across the rankable economies of a technology-stage pair,
  $\mathrm{cov}(\omega_{kits}c_{kits},\EO_{its})/\mathrm{var}(\EO_{its})$, averaged over the eleven
  pairs of the stage. Bar colour is the stage; the bars of one colour sum to 100 per cent within
  a panel. Categories marked * (Technological Readiness, Energy Prices) are constant across
  economies in every pair and so contribute nothing to the right panel; Technology Demand is
  constant except for its BNEF regional and EV-market composites. Categories in the order of
  Table~\ref{tab:eo_structure}. Sample: the 4{,}301 rankable cells of the eleven headline
  technologies. Source: \texttt{economic\_opportunity\_category\_contributions} in
  \texttt{data/processed/outputs/index\_outputs.rds}; \texttt{scripts/paper/fig\_eo.R}.}
  \label{fig:eo_contrib}
\end{figure}

\begin{figure}[tbp]
  \centering
  \includegraphics[width=\textwidth]{eo_heatmap}
  \caption{Position of the coalition members and China in every technology-stage pair of the
  Economic Opportunity pillar. Rows are the 31 core members and 26 partners of the coalition
  (Section~\ref{sec:coalition}), each tier ordered by its members' mean percentile across their
  rankable cells, and China as a reference; columns are the eleven headline technologies within
  each stage. Colour is the economy's percentile rank of $\EO_{its}$ among the rankable economies
  of the pair (equation~\eqref{eq:pct_rank}, ties averaged; 1 is the top economy, 0 the bottom).
  Percentile ranks rather than levels are shown because technology-stage fixed effects account for
  58 per cent of the variance of the level, which would otherwise dominate the colours (levels are
  in Figure~\ref{fig:eo_distribution} and Table~\ref{tab:eo_leaders}). Gray cells with a dot are
  scored but not ranked: every member is sovereign, so these are cells in which more than half of
  the economy's index weight rests on imputed components ($\iota_{its}>0.5$,
  equation~\eqref{eq:imputed_share}). All 33 pairs have status \emph{ok}. Source:
  \texttt{economic\_opportunity\_index} and \texttt{coalition\_members} in
  \texttt{data/processed/outputs/index\_outputs.rds}; \texttt{scripts/paper/fig\_eo.R}.}
  \label{fig:eo_heatmap}
\end{figure}

\begin{figure}[tbp]
  \centering
  \includegraphics[width=\textwidth]{eo_top}
  \caption{The ten rankable economies with the highest Economic Opportunity score in four
  technology-stage pairs: upstream oil, midstream batteries, midstream solar and downstream
  electric vehicles ($n$ is the number of rankable economies in the pair). Each bar's length is
  $\EO_{its}$, stacked by the exact variable decomposition of equation~\eqref{eq:var_decomp}; here
  colour identifies the category, not the stage. The gray segment sums the technology-level
  constants (the stage-mapped TRL Index, WEO global demand, BCG addressable market, overcapacity,
  world-market concentration and commodity-price volatility, whichever the pair carries) and is the
  same for every economy in a pair; the coloured segments are the categories' country-varying
  components, including those filled by a missing-data rule. Coalition members (core and partners)
  are in bold. Source: \texttt{economic\_opportunity\_index} and
  \texttt{economic\_opportunity\_variable\_contributions} in
  \texttt{data/processed/outputs/index\_outputs.rds}; \texttt{scripts/paper/fig\_eo.R}.}
  \label{fig:eo_top}
\end{figure}
